\section{Planar-Magnetics Integration}
\label{sec:disc-magnetics}

The supplied planar inductor and transformer were integrated into the
converter, including terminal assignment, polarity, mounting, and return
paths. Their proximity to switching nodes also affects parasitics and
sensing. The electromagnetic design came from earlier DTU work.

The given component measurements in Table~\ref{tab:mag-inherited-measurements}
are approximately \SI{33}{\micro\henry} for the inductors,
\SI{16.5}{\micro\henry} transformer leakage,
\SIrange{140}{151}{\micro\henry} magnetizing inductance, and a turns
ratio near 8:1 in this thesis's convention. The installed equivalent
circuit must also account for the final A/B assignment and added series
resonant inductors.

At \SI{1}{\mega\hertz}, leakage, winding capacitance, and connection
inductance affect resonance and commutation, while winding proximity and
interconnect resistance affect loss. Power transfer was demonstrated,
but transformer, pump-inductor, and added-inductor losses were not
separated. Resonant current, installed impedances, and component
temperatures are needed to distinguish magnetic-design limitations from
mounting and interconnection losses.

Secondary-winding heating was observed in earlier tests around \SIrange{3.3}{3.6}{\ampere} at the secondary output. Later testing exceeded this current, reaching \SI{10.01}{\ampere} at \SI{13.69}{\volt} with voltage droop. Winding resistance, current distribution, leakage, and output headroom have not been measured together, so the causes of heating and droop remain unresolved.

V3 retained these magnetic constraints; replacement planar magnetics were
not installed within the available time. Transformer optimization may
support the \SI{100}{\watt} objective.

Feedback on the APEC digest raised the question of whether the converter
could support higher loads, including \SI{240}{\watt} USB-PD Extended
Power Range (EPR). Assessing higher-power operation requires the secondary-region and winding heating, input-bridge heating, and output droop to be investigated.
